\section[Original capped sublevels]{Original capped sublevels and a degree-only ACL shortcut}
\label{sec:op3-capped-sublevels}

\paragraph{Scope.}
Return to the original simple connected unweighted graph, physical seed
$v$, $\gamma=(1-\alpha)/(1+\alpha)$, $\bar\alpha=1-\gamma$, and
$\lambda=\epsappr/2$. The following \statusproved{} drafts use original
degrees. They do not assert a bound for every recursive or accelerated
query of a generalized diffusion solver.

\begin{lemma}[Original capped coercivity; \statusproved]
\label{lem:op3-original-capped-coercivity}
Apply \cref{lem:op3-capped-grounding-vwf} with
$h_i=\bar\alpha d_i$, $f=e_v-\lambda d$, and cap $U=1/\bar\alpha$.
For every feasible $x\geq0$,
\begin{equation}
 E_{\rm cap}(x)\geq\lambda\sum_i x_i-\frac1{2\bar\alpha}.
 \label{eq:op3-original-capped-coercivity}
\end{equation}
In particular, $E_{\rm cap}(x)\leq B_E$ implies
$\sum_i x_i\leq(B_E+1/(2\bar\alpha))/\lambda$.
The nonpositive sublevel lies in a box of radius
$1/(2\bar\alpha\lambda)$ in every coordinate.
\end{lemma}
\begin{proof}
Every edge term is nonnegative. At a nonseed vertex, the capped quadratic
grounding is nonnegative and the load term is $\lambda d_ix_i$,
so $q_i(x_i)\geq\lambda x_i$. At the seed, put
$a=f_v+\lambda=1-\lambda(d_v-1)\leq1$. On $[0,U]$,
\[
 q_v(x)-\lambda x=\bar\alpha d_vx^2/2-ax
 \geq-\frac{\max\{a,0\}^2}{2\bar\alpha d_v}
 \geq-\frac1{2\bar\alpha}.
\]
Beyond the cap, the derivative of $q_v(x)-\lambda x$ is
$d_v-a=(d_v-1)(1+\lambda)\geq0$. Its minimum is therefore already
attained on the quadratic interval. Sum the vertex bounds.
\end{proof}

\begin{proposition}[One degree query for weak diffusion; \statusproved]
\label{prop:op3-seed-only-shortcut}
One original seed-degree query suffices for the following ACL certificates.
If $\epsappr d_v\geq1$, return zero. Otherwise, if
$\gamma\leq\epsappr d_v$, return the physical vector $u=e_v/d_v$,
whose ACL output is $\bar\alpha Du=\bar\alpha e_v$.
Both branches cost $O(1)$ words and operations and read no adjacency entry.
If neither branch applies, then $\gamma>\epsappr d_v\geq\epsappr$.
\end{proposition}
\begin{proof}
For zero, the residual is $e_v$, which obeys
$0\leq e_v\leq\epsappr d$ precisely under the stated condition.
For $u=e_v/d_v$ and $M=D-\gamma A$, the residual at $v$ is zero.
Every neighbor has residual $\gamma/d_v\leq\epsappr\leq\epsappr d_i$;
all other residuals are zero. This proves the original ACL certificate
without revealing who the neighbors are. Degree, parameter arithmetic
and the single output coordinate take constant work. The final strict
inequality follows from failure of the two tests and $d_v\geq1$.
\end{proof}

The shortcut is an approximate ACL output, not an exact obstacle solution.
It removes a possible logarithmic dependence on $1/\gamma$ from the input
of a remaining algorithm: on its remaining branch,
$\log(1/\gamma)<\log(1/\epsappr)$. It does not remove polynomial
$1/\bar\alpha$ work from any solver or establish OP3 by itself.

\paragraph{A conditional coercivity transfer.}
If a compressed graph objective obeys the global embedding in
\cref{cor:op3-global-vwf-energy-embedding}, and its input satisfies
$E(x)\geq\lambda\sum_i x_i-C_0$ on $x\geq0$, then
\[
 \widehat E(x)\geq\lambda\sum_i x_i-2C_0-2\Xi.
\]
This follows by substituting $x/2$ in the input bound. It gives a concrete
way to track sublevels through one compression. To use repeated transfers,
one must bound the actual compression depth and total accumulated errors;
counting only a single invocation would not suffice.

\paragraph{Audit.}
\texttt{capped\_sublevel\_bounds.py} checks all 30 connected atlas graphs
on two through five vertices, every physical seed, three teleportation
values and four original ACL tolerances: 1,644 cases. Exact quadratic
minima and final-ray slopes certify all 7,716 vertex bounds; 16,440
feasible vectors check the original capped objective, including values
beyond the cap. There are 411 zero shortcuts, 685 seed-only shortcuts and
548 continuations, all checked against the original matrix. Implicit
center-seeded stars with $16$, $10^{12}$ and $10^{30}$ leaves certify one
degree query, zero adjacency reads and the exact threshold tie. The graph
enumeration and dense obstacle solutions are validators, not charged
operations of the constant-work shortcut.
